% Options for packages loaded elsewhere
\PassOptionsToPackage{unicode}{hyperref}
\PassOptionsToPackage{hyphens}{url}
\documentclass[
]{article}
\usepackage{xcolor}
\usepackage{amsmath,amssymb}
\setcounter{secnumdepth}{-\maxdimen} % remove section numbering
\usepackage{iftex}
\ifPDFTeX
  \usepackage[T1]{fontenc}
  \usepackage[utf8]{inputenc}
  \usepackage{textcomp} % provide euro and other symbols
\else % if luatex or xetex
  \usepackage{unicode-math} % this also loads fontspec
  \defaultfontfeatures{Scale=MatchLowercase}
  \defaultfontfeatures[\rmfamily]{Ligatures=TeX,Scale=1}
\fi
\usepackage{lmodern}
\ifPDFTeX\else
  % xetex/luatex font selection
\fi
% Use upquote if available, for straight quotes in verbatim environments
\IfFileExists{upquote.sty}{\usepackage{upquote}}{}
\IfFileExists{microtype.sty}{% use microtype if available
  \usepackage[]{microtype}
  \UseMicrotypeSet[protrusion]{basicmath} % disable protrusion for tt fonts
}{}
\makeatletter
\@ifundefined{KOMAClassName}{% if non-KOMA class
  \IfFileExists{parskip.sty}{%
    \usepackage{parskip}
  }{% else
    \setlength{\parindent}{0pt}
    \setlength{\parskip}{6pt plus 2pt minus 1pt}}
}{% if KOMA class
  \KOMAoptions{parskip=half}}
\makeatother
\setlength{\emergencystretch}{3em} % prevent overfull lines
\providecommand{\tightlist}{%
  \setlength{\itemsep}{0pt}\setlength{\parskip}{0pt}}
\usepackage{bookmark}
\IfFileExists{xurl.sty}{\usepackage{xurl}}{} % add URL line breaks if available
\urlstyle{same}
\hypersetup{
  hidelinks,
  pdfcreator={LaTeX via pandoc}}

\author{}
\date{}

\begin{document}

\textbf{NurseSync AI}

\textbf{Intelligent Audio-Based Clinical Handoff System}

Project Proposal for UCS503P

Submitted to: Dr. Asif

\textbf{Thapar Institute of Engineering and Technology}

Shivansh Sharma( Roll No: 1024030732 )

Ishaan Bhalla( Roll No: 1024030950 )

Pushp Batheja( Roll No: 1024030941 )

August 15, 2026

\section{1 Higher-order goal}\label{higher-order-goal}

This project aims to improve patient safety and continuity of care by
reducing the time, inconsistency, and error rate involved in
nurse-to-nurse clinical handoffs. While the underlying goal --- safer,
more reliable communication of patient information --- is broader than
any single application, the immediate engineering objective is to
translate structured, accurate handoff documentation into a measurable,
deployable software solution: a system that lets a nurse speak a handoff
and converts that speech into structured, searchable, and auditable
clinical records.

\section{2 Time-to-value}\label{time-to-value}

\begin{itemize}
\item
  We will deliver meaningful increments quickly by building the core
  patient/handoff workflow (login, patient list, patient profile,
  structured handoff storage) before adding any AI/ML capability.
\item
  We will prioritize fast feedback over perfection by following a staged
  roadmap --- development environment, backend, database, patient
  system, mobile frontend, audio capture, speech-to-text, NLP
  extraction, summary/timeline, RAG assistant, ML risk model, testing,
  deployment --- validating each stage before moving to the next.
\item
  Success in the first iteration is defined narrowly: a nurse can log
  in, view a patient, and create and retrieve a structured (even if
  manually entered) handoff record. AI/ML layers are added incrementally
  on top of this working core.
\end{itemize}

\section{3 Problem Statement}\label{problem-statement}

Nurses routinely hand off patient information verbally or in free-text
notes at the end of a shift --- for example, describing a
patient\textquotesingle s diagnosis, vitals, medications administered,
and pending tests in a single spoken narrative. Reports indicate nurses
can spend a substantial share of a shift on manual charting, time taken
directly from patient care. Common pain points include:

\begin{itemize}
\item
  \textbf{Fragmentation:} handoff information is typed manually, spread
  across multiple notes, or communicated purely verbally, causing
  details to be lost or inconsistently recorded.
\item
  \textbf{Inconsistent structure:} different nurses document
  differently, making handoffs harder to scan quickly during a busy
  shift change.
\item
  \textbf{Poor searchability:} once written, information is difficult to
  search or retrieve later (e.g., "what was this
  patient\textquotesingle s oxygen level this morning?").
\item
  \textbf{Handoff errors:} incomplete or misinterpreted information
  during nurse-to-nurse handoffs is a recognized contributor to
  preventable medical errors and missed critical details.
\item
  \textbf{Accountability gaps:} without reliable, time-stamped records,
  it is difficult to reconstruct exactly what happened and when during a
  patient\textquotesingle s care.
\end{itemize}

Why this matters now (contextual relevance): On a busy ward, even small
improvements in how quickly and accurately information is captured and
transferred can meaningfully reduce missed or delayed care.

\section{}\label{section}

\section{4 Proposed Solution}\label{proposed-solution}

\subsection{4.1 Overview}\label{overview}

Build a mobile-first "speak-to-structured-handoff" system with the
following components:

\begin{itemize}
\item
  \textbf{Voice handoff capture:} a nurse records a spoken handoff for a
  specific patient instead of typing it.
\item
  \textbf{Speech-to-text:} the recorded audio is transcribed to a text
  transcript.
\item
  \textbf{Medical entity extraction (NLP):} the transcript is parsed
  into structured fields --- diagnosis, medications, vitals, allergies,
  tests, and events.
\item
  \textbf{Structured SBAR handoff + AI summary:} extracted information
  is organized into a standard
  Situation/Background/Assessment/Recommendation format and a short,
  scannable summary.
\item
  \textbf{Timeline generator:} every extracted event (medication given,
  vitals recorded, test ordered) becomes a timestamped entry in a
  longitudinal patient timeline.
\item
  \textbf{Missing-information detection:} the system checks a handoff
  against an expected field set (diagnosis, vitals, allergies, pending
  tests, recommendation) and flags gaps.
\item
  \textbf{RAG-grounded AI assistant:} nurses can ask natural-language
  questions about a patient ("what changed today?") and receive answers
  grounded strictly in that patient\textquotesingle s stored records.
\item
  \textbf{Experimental risk model:} a machine-learning classifier
  estimates deterioration risk (low/medium/high) from structured vitals,
  with an explanation of contributing factors.
\end{itemize}

\subsection{4.2 Core workflow}\label{core-workflow}

\begin{enumerate}
\def\labelenumi{\arabic{enumi}.}
\item
  Nurse selects a patient and taps "Record Handoff."
\item
  Nurse speaks the handoff naturally; audio is recorded and uploaded.
\item
  Backend transcribes the audio, extracts medical entities and events,
  and generates an SBAR handoff, a concise summary, and timeline
  entries.
\item
  The completeness checker flags any missing required fields, and the
  risk model (where applicable) scores deterioration risk from
  structured vitals.
\item
  Everything is stored; the next nurse opens the patient and sees the
  summary, timeline, risk indicator, full handoff history, and can query
  the AI assistant.
\end{enumerate}

\subsection{4.3 Operational constraints and safety
boundaries}\label{operational-constraints-and-safety-boundaries}

\begin{itemize}
\item
  Keep the mobile app usable hands-free at the bedside, on standard
  hospital devices.
\item
  The system provides AI-assisted organization, summarization,
  retrieval, and experimental risk prediction to support the handoff
  workflow --- it does not diagnose patients or make autonomous clinical
  decisions. The nurse remains responsible for verifying all
  information.
\item
  RAG answers must be grounded in retrieved patient records; the
  assistant should not generate unsupported clinical claims.
\item
  For this student prototype, use synthetic or de-identified patient
  data rather than real patient recordings.
\end{itemize}

\section{}\label{section-1}

\section{5 Solution Approach}\label{solution-approach}

This is a software engineering project with an applied-AI component, so
the emphasis is on integrating mature, existing AI/ML components into a
correct, well-tested pipeline --- rather than researching new model
architectures.

\subsection{5.1 AI/ML component integration
\ldots non-research}\label{aiml-component-integration-non-research}

\begin{itemize}
\item
  \textbf{Speech-to-text:} use an existing pretrained model (e.g.,
  Whisper) rather than training a speech model from scratch.
\item
  \textbf{Entity extraction \& summarization:} use an LLM/NLP pipeline
  via prompting or lightweight fine-tuning to pull structured fields out
  of transcripts and produce the SBAR/summary text.
\item
  \textbf{RAG question answering:} chunk and embed patient records,
  retrieve the most relevant records for a given question, and pass only
  that retrieved context to the LLM so answers stay grounded in stored
  data instead of being generated freely.
\item
  \textbf{Missing-information detection:} start as a rules + NLP check
  against an expected field checklist; consider a trained classifier
  later only if a suitable labeled dataset becomes available.
\item
  \textbf{Risk classification:} compare established, well-understood
  models --- Logistic Regression, Random Forest, XGBoost --- rather than
  assuming one is best, and report feature-level contributions
  (explainability) alongside each prediction.
\end{itemize}

\subsection{\texorpdfstring{5.2 Mobile / backend architecture
}{5.2 Mobile / backend architecture }}\label{mobile-backend-architecture}

\begin{itemize}
\item
  \textbf{Frontend:} a Flutter mobile app covering login, dashboard,
  patient list, patient profile, audio recorder, timeline view, and AI
  assistant chat.
\item
  \textbf{Backend API:} FastAPI service exposing endpoints such as
  /login, /patients, /handoffs, /handoffs/\{id\}/process,
  /patients/\{id\}/timeline, /patients/\{id\}/ask, and
  /patients/\{id\}/risk.
\item
  \textbf{Database:} PostgreSQL storing Patient, Handoff, Vitals,
  Medication, Timeline, and Risk tables with clear relations back to
  each patient.
\end{itemize}

\subsection{5.3 CI/CD and fast delivery
enabler}\label{cicd-and-fast-delivery-enabler}

\begin{itemize}
\item
  Automatically build and run backend unit tests and linting on every
  change.
\item
  Produce repeatable deployment artifacts to a staging environment.
\item
  Enable safe, frequent releases of each incremental stage in the
  roadmap (backend → database → patient system → mobile app → audio →
  speech AI → NLP → summary/timeline → RAG → ML → testing → deployment).
\end{itemize}

\section{6 Evaluation Criterion}\label{evaluation-criterion}

We define success using measurable metrics tied to the core workflow.

\subsection{6.1 Primary evaluation
metric}\label{primary-evaluation-metric}

Handoff Documentation Time (HDT): time from the start of voice recording
to a fully structured, stored handoff record (transcript + SBAR +
summary + timeline entries).

\begin{itemize}
\item
  Target: median HDT well below the time required to manually type an
  equivalent handoff, measured in a small user study.
\item
  Attribution: measured from timestamps logged at recording start and at
  the handoff-stored event in the database.
\end{itemize}

\subsection{6.2 Secondary evaluation
metrics}\label{secondary-evaluation-metrics}

\begin{itemize}
\item
  \textbf{Entity extraction accuracy:} precision/recall/F1 of extracted
  fields (diagnosis, medication, vitals, allergy, tests) against
  manually labeled transcripts.
\item
  \textbf{Completeness-detection accuracy:} precision/recall of the
  missing-information checker against handoffs with intentionally
  omitted fields.
\item
  \textbf{RAG answer groundedness:} proportion of AI assistant answers
  that are fully supported by the retrieved patient records, reviewed
  manually on a test set of questions.
\item
  \textbf{Risk model performance:} accuracy, precision, recall, F1, and
  ROC-AUC, with particular attention to recall/sensitivity and false
  negatives given the medical stakes.
\item
  \textbf{Nurse-reported usability:} clarity and usefulness of
  summaries/timeline using a simple 1--5 feedback question.
\item
  \textbf{Reliability:} availability of the staging/production API
  during testing sessions.
\end{itemize}

\subsection{6.3 Pilot validation plan
}\label{pilot-validation-plan-high-level}

\begin{itemize}
\item
  Recruit a small group of nursing students or clinically-informed
  volunteers (e.g., 5--10) to record handoffs for synthetic patient
  scenarios.
\item
  Compare structured-handoff creation time against a manual-typing
  baseline for the same scenarios.
\item
  Collect qualitative feedback on summary/timeline usefulness,
  missing-information flags, and AI assistant answer quality.
\end{itemize}

\section{7 Scalability}\label{scalability}

\begin{enumerate}
\def\labelenumi{\arabic{enumi}.}
\item
  \textbf{Scalable (theoretically):} the system should support growth in
  number of patients, handoffs, and concurrent app usage by:

  \begin{itemize}
  \item
    decoupling the Flutter frontend from the FastAPI backend,
  \item
    using a stateless REST API design,
  \item
    indexing and paginating common queries (patient lists, timelines),
  \item
    processing audio transcription asynchronously (queue-based) so
    long-running AI jobs don\textquotesingle t block the API.
  \end{itemize}
\item
  \textbf{Operationally deployable (practically):} the system should run
  on common infrastructure:

  \begin{itemize}
  \item
    a managed PostgreSQL database,
  \item
    containerized backend deployment,
  \item
    a CI/CD pipeline for staging/production,
  \item
    a vector store for RAG retrieval that can scale independently of the
    relational database.
  \end{itemize}
\end{enumerate}

\section{8 Engine availability
heuristic}\label{engine-availability-heuristic}

A mature set of "engines" (tooling/services) is available to implement
this project without inventing new infrastructure or model
architectures:

\begin{itemize}
\item
  Speech-to-text engine (e.g., Whisper) for audio → transcript.
\item
  Hosted or open-source LLM for entity extraction, SBAR structuring,
  summarization, and RAG-based answer generation.
\item
  Embedding + vector search library for retrieval-augmented generation.
\item
  FastAPI as the REST framework; Flutter as the cross-platform mobile
  framework.
\item
  Managed PostgreSQL for relational storage.
\item
  Token-based authentication (e.g., JWT).
\item
  Object storage for audio files.
\item
  scikit-learn / XGBoost for risk classification experiments.
\item
  pytest for backend unit/integration tests; a CI runner and staging
  deployment target.
\end{itemize}

We will use these mature components to focus on pipeline design, data
structuring, correctness, and clinical-safety framing rather than
building new infrastructure or novel model architectures.

\section{9 Project scope and
deliverables}\label{project-scope-and-deliverables}

\subsection{9.1 Initial deliverable \ldots first
iteration}\label{initial-deliverable-first-iteration}

\begin{itemize}
\item
  Login/authentication and a Flutter app shell connected to the FastAPI
  backend.
\item
  Patient list and patient profile screens backed by a PostgreSQL schema
  (Patient, Handoff, Vitals, Medication, Timeline, Risk).
\item
  Manual handoff entry and retrieval (structured fields, no AI yet), to
  validate the core data model end to end.
\item
  CI: build + backend unit tests + linting.
\item
  CD to staging.
\end{itemize}

\subsection{9.2 Subsequent deliverables}\label{subsequent-deliverables}

\begin{itemize}
\item
  Audio recording, upload, and storage.
\item
  Whisper-based speech-to-text integration.
\item
  Medical NLP entity extraction feeding an auto-generated SBAR handoff
  and AI summary.
\item
  Automatic timeline generation from extracted events.
\item
  Missing-information / completeness checker.
\item
  RAG-based AI assistant for patient-specific question answering.
\item
  Experimental ML deterioration risk model with explainability, plus
  expanded testing and deployment.
\end{itemize}

\section{}\label{section-2}

\section{10 Risks and mitigations}\label{risks-and-mitigations}

\begin{itemize}
\item
  \textbf{Patient data privacy/consent:} use only synthetic or
  de-identified data for the prototype; no real patient recordings; plan
  for stronger safeguards (access control, audit logs) before any real
  deployment.
\item
  \textbf{Overclaiming AI capability:} explicitly label the risk model
  and summaries as decision-support, not diagnosis; keep the nurse
  responsible for verification; constrain RAG answers to retrieved
  records to avoid hallucination.
\item
  \textbf{Speech recognition accuracy in clinical/noisy settings:}
  evaluate transcription quality on medical vocabulary; give the nurse a
  quick review/correction step before a handoff is finalized.
\item
  \textbf{False negatives in risk prediction:} prioritize
  recall/sensitivity over raw accuracy; tune thresholds; keep a human in
  the loop for any high-risk flag.
\item
  \textbf{Staff adoption:} keep the voice-recording workflow faster than
  typing; make structured output easy to review and edit; minimize extra
  steps versus current practice.
\item
  \textbf{Scope creep across many AI/ML stages:} follow the staged
  roadmap and ship each stage (backend → database → frontend → audio →
  speech AI → NLP → summary/timeline → RAG → ML → testing → deployment)
  as a working increment before starting the next.
\end{itemize}

\section{11 Summary}\label{summary}

NurseSync AI proposes a deployable, voice-driven clinical handoff system
that converts spoken nurse-to-nurse handoffs into structured,
timestamped, and searchable clinical records, augmented with AI
summarization, automatic timeline generation, missing-information
detection, and a RAG-grounded assistant for patient-specific question
answering. An experimental, explainable ML deterioration-risk model adds
a decision-support layer without replacing clinical judgment. The
project emphasizes engineering integration of mature, existing AI/ML
components, incremental delivery through CI/CD, and measurable
evaluation --- handoff documentation time, extraction accuracy, and
risk-model recall --- while maintaining clear safety boundaries: the
system organizes, summarizes, and retrieves information to support the
nurse, and does not diagnose or make autonomous care decisions.

\end{document}
